\item \textbf{Fisión espontánea del Uranio-236 en fragmentos de Bromo y Lantano.}
Un núcleo de $^{236}_{92}\text{U}$ inicialmente en reposo se fisiona espontáneamente en dos fragmentos: $^{87}_{35}\text{Br}$ y $^{149}_{57}\text{La}$.
\begin{enumerate}
    \item Demuestre usando conservación del momento lineal que la fracción de la energía cinética total llevada por el fragmento $m_1$ es $K_1/K_{\text{tot}} = m_2 / (m_1 + m_2)$.
    \item Calcule la energía total de desintegración $Q$ liberada en el proceso.
    \item Determine cómo se reparte cuantitativamente la energía cinética entre ambos fragmentos primarios.
    \item Calcule las velocidades individuales de cada fragmento inmediatamente después de la fisión.
\end{enumerate}